\documentclass[11pt]{article}
\usepackage[letterpaper,margin=1.05in]{geometry}
\usepackage[T1]{fontenc}
\usepackage{lmodern,amsmath,amssymb,amsthm,microtype,booktabs}
\usepackage[round,authoryear]{natbib}
\usepackage[colorlinks=true,linkcolor=blue,citecolor=blue,urlcolor=blue]{hyperref}
\newtheorem{proposition}{Proposition}
\newtheorem{definition}[proposition]{Definition}
\theoremstyle{definition}
\newtheorem{remark}[proposition]{Remark}
\newcommand{\W}{\mathcal W}
\newcommand{\U}{\mathcal U}
\title{Local presents and unrestricted existence:\\ the coverage problem for dynamical presentism}
\author{Working manuscript}
\date{21 October 2026}
\begin{document}
\maketitle

\begin{abstract}
Brading's law-constitutive present is system-relative: a present is the
spatiotemporal region needed to sustain a dynamical unity. This paper asks
whether that local present restricts unrestricted concrete existence, as
presentism requires. It does not, by itself. Local boundedness of each
system-present is compatible with the union of all presents covering the whole
history; the coverage objection is exact. Brading's own reformulation and her
interaction passage support a system-relative existence reading. Recovering a
global restriction requires an actuality or privilege premise, restrictions on
coexistence, or a revised presentism, each with a cost not settled by physics.
The result is a limited verdict: dynamical individuation is insufficient for the
presentist existential claim. The paper does not refute Brading or establish
eternalism.
\end{abstract}

\section{Introduction}

Brading's \emph{Presentism as an Empirical Hypothesis} offers one of the most
careful attempts to keep a presentist ontology alive within special relativity
\citep{Brading2013}. She rejects the shared assumption that spacetime is the
ontological principle of unity of what exists, and replaces it with a
law-constitutive principle: the dynamical laws constitute genuine unities, and
the present of a system is the spatiotemporal region needed to sustain its
dynamics \citep[pp.~1104--1109]{Brading2013}. The present is therefore local and
system-relative, not a global simultaneity slice.

This paper asks a narrower question than whether presentism or eternalism is
true:

\begin{quote}
Under what conditions does a physically defined local present constitute a
restriction on what exists, rather than merely a division of an eternalist
history into dynamical unities?
\end{quote}

The answer defended here is negative and limited. Brading's local present does
not, by itself, establish a restriction on unrestricted concrete existence. Her
own reformulation,

\begin{quote}
(P1*) For each and every thing, that thing exists only presently, where it is
something to be determined empirically \citep[p.~1110]{Brading2013},
\end{quote}

is a per-thing, system-relative existence claim, not the original
``all and only things that exist now are real'' \citep[p.~1101]{Brading2013}.
Local boundedness of each system-present is compatible with the union of all
presents covering the entire history. Recovering a global restriction requires
an additional actuality or privilege premise, a restriction on which presents
coexist, or a revision of what presentism claims. Each repair has a cost that
is not settled by the physics.

The contribution is not a new general theory of time. It is a focused
reconstruction of one existence bridge and a precise objection to it. The paper
does not refute Brading, does not prove eternalism, and does not claim that
every presentism fails.

\section{The existential contrast}

\subsection{Two premises and a reformulation}

Brading begins from the familiar tension between

\begin{itemize}
\item (P1) All and only things that exist now are real;
\item (P2) Special relativity is a complete account of spatiotemporal
structure \citep[p.~1101]{Brading2013}.
\end{itemize}

P1 is a global restriction on existence: what there is is what is present, and
events outside the present are not real. P2 says that nothing is missing from
the special-relativistic account of spatiotemporal structure. The apparent
conflict is that relativity supplies no global present.

Brading's response is not to add a preferred foliation. It is to reject the
assumption that spacetime is the ontological ground of unity, and to ground
unity in the dynamical laws instead \citep[pp.~1104--1108]{Brading2013}.
Spacetime becomes an epistemic principle of unity, the framework in which we
organize knowledge beyond the here and now. The present becomes the region
needed to sustain a dynamical system, and P1 is rewritten as P1*.

P1* is the textual pivot of this paper. It quantifies over things and assigns
each a present. It does not say that there is a global domain of existence
smaller than the world history. That difference is the subject of the coverage
objection.

\subsection{Existence, presence, unity and endurance}

Four notions must be kept apart:

\begin{enumerate}
\item \emph{Existence}: the unrestricted claim that an event is part of what
there is.
\item \emph{Presence}: an event's belonging to a system's sustaining region.
\item \emph{Unity}: the law-constituted integration that makes a system one
thing.
\item \emph{Endurance}: the persistence of that unity through change.
\end{enumerate}

Brading's argument moves from unity to presence. The further move from presence
to a restriction on existence is the one at issue. Confusing presence with
existence makes the coverage problem invisible: every event can be present in
some system's present while every event exists.

\subsection{The target contrast}

Let $\W$ be the set of concrete events in the world history, and for each
genuine dynamical system $i$ let $\U_i\subseteq\W$ be its system-sustaining
present.

\begin{definition}[Global restriction]
A reading imposes a \emph{global existential restriction} if there is a fact
about what exists simpliciter and $\W$ contains events that do not satisfy it.
\end{definition}

\begin{definition}[Local boundedness]
A reading is \emph{locally bounded} if $\U_i$ is a proper subregion of $\W$ for
every system $i$.
\end{definition}

Presentism, in the sense of P1, requires a global restriction. Eternalism denies
one. Brading's system-relative present is locally bounded; whether it imposes a
global restriction is exactly what the coverage objection puts in question.

\section{Brading's law-constitutive present}

Brading's positive proposal has three parts.

First, the laws constitute genuine unities. A physical thing, simple or
composite, is what it is in part because it satisfies the laws; the laws also
constitute the parts of a composite and their persistence through change
\citep[pp.~1105--1108]{Brading2013}. This is offered as an alternative to
essentialism and haecceitism, and as an account of numerical identity through
qualitative change.

Second, spacetime is demoted from an ontological to an epistemic principle of
unity. It is the framework in which we organize knowledge of the not-here and
not-now, not the ground of what exists \citep[pp.~1108--1109]{Brading2013}.

Third, the present is the region needed to sustain the dynamical system:

\begin{quote}
The present is a spatiotemporal region of whatever size is necessary to sustain
the dynamical system in question \citep[p.~1109]{Brading2013}.
\end{quote}

The region may be much smaller than the block, or it may be the size of the
block. The 2015 follow-up develops the same system-sustaining-region
formulation \citep{Brading2015}. If the only genuine system is the entire universe, Brading says that
``this version of presentism is defeated and the block universe triumphs''
\citep[p.~1109]{Brading2013}. The empirical question is whether there are
genuine subsystems whose sustaining regions are proper parts of the whole.

This is an elegant proposal. It avoids pluralistic solipsism by denying that
spacetime is the ground of unity, and it makes the extent of the present an
empirical matter. The question is what it claims about existence.

A distinguished predecessor is Stein's event-relative notion of definiteness:
the state at a point is definite as of that point, and the relation is local
rather than a global present \citep{Stein1991}. Stein's printed theorem
(p.~149) lists reflexivity, transitivity, invariance under
time-orientation-preserving automorphisms and one past-pointing related pair.
It is not self-contained: the proof also uses the non-universality condition
(iv) and the prose premise that every event in the causal past---timelike or
null---is definite as of the point. The reflexive strict chronological past is
non-universal but omits past-null pairs, so condition (iv) alone is not enough.
The local-present strategy inherits this distinction; it does not inherit a
global restriction.

\section{The coverage objection}

\subsection{Local boundedness is not global restriction}

Local boundedness says, for every system $i$,

\[
\U_i\subsetneq\W .
\]

It does not imply

\[
\bigcup_i \U_i\subsetneq\W .
\]

If every event belongs to at least one system-present, the family
$\{\U_i\}$ covers $\W$. The account then divides the history into dynamical
unities; it does not exclude any event from existence. The present can be
``local, not global'' \citep[p.~1110]{Brading2013} while the union of local
presents is the whole history.

\begin{proposition}[Coverage]
Local boundedness of every $\U_i$ is consistent with
$\bigcup_i\U_i=\W$.
\end{proposition}
\begin{proof}
Let $\W$ be a nonempty set and let $\{\U_i\}$ be any cover of $\W$ by proper
subsets. For example, take two overlapping proper subsets whose union is $\W$.
Every $\U_i$ is locally bounded, and the union is $\W$.
\end{proof}

The proposition is elementary. Its philosophical point is not the set theory
but the missing premise: local boundedness does not supply a global
restriction. Overlap does not help, because overlapping presents make coverage
easier. Disjointness does not help either, because a partition of $\W$ into
disjoint proper parts still covers $\W$.

\subsection{Brading's text points to system-relative existence}

The coverage objection is not an external imposition. Brading's own gloss
supports the system-relative reading.

P1* quantifies over things and assigns each a present. It is not P1. More
strikingly, the interaction passage says:

\begin{quote}
The ground of what is real is the dynamics, and we belong to the same world as
whatever we interact with, and the rest of that world is as real as we are.
\ldots\ Does my son exist relative to me when he is in London and I am here? Of
course. There are plenty of interactions going on that link us. Is there a
determinate fact of the matter about what he is doing right now? No.
\citep[p.~1110]{Brading2013}
\end{quote}

This is a relational existence claim: the son exists relative to the speaker
because of interaction; there is no determinate global fact about his present.
The rest of the world is ``as real as we are.'' That is system-relative
existence, not the global restriction of P1.

Brading's empirical defeat condition points the same way: if the only genuine
system is the whole universe, the view is defeated and the block triumphs
\citep[p.~1109]{Brading2013}. This is a conditional about which systems exist,
not a global actuality principle.

Whether the view should still be called ``presentism'' is partly verbal;
Brading herself worries that ``labeling it `presentism' dooms the view''
\citep[p.~1101]{Brading2013}. The substantive point is that the global
restriction is not established.

\section{Four repairs and their costs}

If one wants the original global restriction, at least one repair is needed.

\paragraph{1. System-relative existence.}
Replace a single unrestricted existence predicate with ``exists-for-system-$i$.''
This blocks the coverage objection by changing the question, but it gives up
the global presentist claim: there is no system-independent fact that future
events do not exist. The view becomes a relativized ontology. Brading's
interaction passage is naturally read this way.

\paragraph{2. Privileged actual systems.}
Only some systems' presents are actual. This restores a restriction, but the
selection principle is not derived from the physics or from dynamical unity. It
needs an independent actuality premise; without it the choice is arbitrary, and
with it the view may reintroduce a privileged perspective.

\paragraph{3. Restrictions on coexistence.}
Not all local presents can coexist: a maximal consistent set, a single actual
present, or a dynamics of present-selection. This risks reintroducing a global
present or a privileged foliation, and it needs a rule the physics does not
supply.

\paragraph{4. Revised presentism.}
Give up global exclusion and accept local or functional becoming. This is
coherent, but it is a different thesis: it may be an eternalist history with a
local becoming, not presentism as the restriction on existence.

The paper does not select among these. It argues that at least one is needed,
and that each has a cost not settled by the physics. That is the limited
verdict.

\section{What the argument owes}

An objection to a bridge should specify what the bridge needs, not demand a
complete rival dynamics. The relevant standard here is admissibility and
explanatory comparison, not a likelihood model of every ontology.

For an admissibility claim, one must state:

\begin{enumerate}
\item the ontology's constitutive commitments;
\item the permitted auxiliary changes;
\item the eligibility evidence;
\item the bridge from the physical structure to the existential claim;
\item what would defeat the claim.
\end{enumerate}

Brading's commitments are law-constituted unity, a system-relative present and
P1*. The eligibility condition is that genuine subsystems exist and their
sustaining regions are proper parts. The bridge is from presence to existence.
The defeat condition is that the universe is the only system. The coverage
objection shows that the bridge is missing. That is enough to make the
criticism bite without demanding a full likelihood model of eternalism.

This is not a licence for vagueness. Background evidence that fixes the
eligible packages must be distinguished from the evidence that compares them.
Filtering out unsuccessful realizations changes the evidential base, and the
cost of the filter must be reported. A criterion can be prejudicial before any
result is inspected; preregistration is not justification. A principled
revision must state what changed, why the change is licensed by the ontology's
commitments, what else it affects, whether the revised package still deserves
the name, and what would have counted against it.

\section{Favoring without exclusion}

The conclusion is not that the evidence is silent. A package can be favored
without excluding every rival.

Treat an ontology as a composite hypothesis: a class of admissible
realizations, a prior over them, and a marginal likelihood. If one ontology has
one realization with $p(D\mid O_1)=1$ and a rival has ten equally weighted
realizations, one with $p(D\mid r)=1$ and nine with $0$, then
$p(D\mid O_2)=0.1$ and the Bayes factor is ten. The rival retains a compatible
realization; evidential parity does not survive.

The prior must be justified by earlier physical evidence, explanatory
achievements, symmetry or an explicitly stated ignorance prior. If it is not,
the verdict is conditional on a stipulated prior, not evidence simpliciter.

For Brading, the presentist composite includes realizations that differ in
which systems are genuine and how presence relates to existence. The
system-relative realization may be behaviorally indistinguishable from an
eternalist history with local presents; if so, the comparison does not favor
presentism. A privileged-actuality realization may favor it, but the privilege
premise is doing the work. That is a substantive comparative verdict, not a
catalogue of failed exclusions.

\section{Strongest objections}

\paragraph{Saunders: the equivalence-relation requirement.}
Saunders argues that presentism needs an invariant, definable, non-trivial
equivalence relation of co-presence \citep{Saunders2000}. The definability
results leave no such relation: the surviving relations are past-chronological,
past-causal or universal, and none is a symmetric equivalence relation of a
present except the trivial universal one. A system-relative present is not an
equivalence relation among all events; it is a region assigned to a system.
Brading's view therefore either gives up the equivalence-relation requirement
and with it the standard presentist claim, or needs an additional structure
that defines co-presence. The coverage objection is the set-theoretic shadow of
the first horn.

\paragraph{Pooley: taking both passage and relativity seriously.}
Pooley distinguishes two easy routes: a deflationary analysis of passage
compatible with the block, which does not take passage seriously, and a
privileged global present, which does not take relativity seriously
\citep{Pooley2013}. His hard route eschews a global Now and links passage to
the open future. Brading's local present is a candidate for that route, but the
coverage objection shows that it can be read as a division of the block into
system-relative presents, with the rest of the world ``as real as we are.''
That is the deflationary horn. If the view adds an actuality or privilege
premise to get a real restriction, it moves toward the second horn or a hidden
privileged structure.

The paper does not refute either challenge. It concedes that it does not
vindicate robust passage. Its contribution is to show that the local-present
route does not by itself restore the presentist existential restriction. That
is a negative result about one bridge, not a proof of eternalism.

\section{Conclusion}

Brading's law-constitutive present is a serious alternative to both the block
and the privileged-now strategy. Its system-relative present is locally
bounded, and its law-constitutive unity is a genuine explanatory proposal. But
local boundedness does not imply a global existential restriction. The coverage
objection is exact: the union of system-presents can be the whole history.
Brading's own P1* and interaction passage support a system-relative existence
reading, not P1. Recovering P1 requires an actuality or privilege premise,
restrictions on coexistence, or a revised presentism; each has a cost not
settled by the physics.

The limited verdict is therefore: dynamical individuation alone is insufficient
for the presentist existential claim. This does not refute Brading, does not
establish eternalism, and does not claim that every presentism fails. It
identifies the missing premise and the price of supplying it.

\paragraph{What would change the verdict.}
A proof that the actual system-presents cannot cover $\W$ would restore the
restriction without an actuality premise. A physically motivated principle
selecting privileged actual systems would supply the missing premise, though
its cost would still need defense. A revised P1* that explicitly restricts the
domain of existence, rather than assigning each thing a present, would change
the target. A demonstration that system-relative existence is the same global
claim under another notation would defeat the system-relative reading.

\paragraph{Limits.}
This is a philosophical argument, not an empirical prediction or an ontology
selection. The coverage objection is elementary set theory; the contribution is
the focused reconstruction and cost analysis. Brading 2013 is read at full text;
Brading 2015 is used at abstract/indexed depth for the system-sustaining-region
formulation. Saunders and Pooley are read at the preprint depths stated in the
accompanying note. No global priority or novelty claim is made.

\bibliographystyle{plainnat}
\bibliography{references}
\end{document}
